\documentclass[11pt,a4paper]{article}
% Estilo comum dos PDFs individuais da Sprint 1.
\usepackage[utf8]{inputenc}
\usepackage[T1]{fontenc}
\usepackage[brazilian]{babel}
\usepackage[a4paper,margin=2.35cm]{geometry}
\usepackage{xcolor}
\usepackage{hyperref}
\usepackage{enumitem}
\usepackage{listings}
\usepackage{parskip}
\usepackage{needspace}
\definecolor{azul}{RGB}{20,63,112}
\hypersetup{colorlinks=true,linkcolor=azul,urlcolor=azul}
\setlist[itemize]{leftmargin=*,itemsep=3pt}
\lstset{basicstyle=\ttfamily\small,frame=single,breaklines=true,columns=fullflexible}
\sloppy
\newcommand{\codigo}[1]{\texttt{#1}}
\newcommand{\titulo}[2]{\begin{center}{\LARGE\bfseries #1\par}\vspace{4pt}{\large #2\par}\vspace{8pt}\rule{0.8\textwidth}{0.4pt}\par\vspace{5pt}{\small Matheus Moreli Durão Silva -- Matrícula 12621CBS217\par Cibersegurança - Bacharelado - Noturno\par}\end{center}}
\newcommand{\cardfonte}[1]{\par\smallskip\noindent\textit{Referência: }\href{#1}{Trello - enunciado e checklist do card}.}

\begin{document}\titulo{S1\_18 - Branches e merge}{Entregar uma funcionalidade com rastreabilidade}
\section*{Objetivo do card}
Aqui a branch deixa de ser apenas conceito: uma funcionalidade deve ser criada em linha separada, documentada e integrada. O entregável é a funcionalidade em branch, com seu caminho até o merge registrado.
\section*{Por que não editar tudo na \codigo{main}?}
A \codigo{main} representa uma versão estável. Uma branch como \codigo{feature/resumo-inventario} isola a criação do resumo: é possível testar, revisar diferenças e corrigir antes de afetar a versão principal. A branch não cria uma cópia permanente do projeto; é uma referência a uma sequência de commits que depois pode ser integrada.
\section*{Passo a passo}
\begin{enumerate}
\item Partir de uma \codigo{main} atualizada e criar a branch com \codigo{git switch -c feature/resumo-inventario}.
\item Implementar uma única finalidade, testar e criar um commit descritivo.
\item Enviar a branch, abrir PR e ler a comparação de arquivos.
\item Se houver conflito, decidir o conteúdo correto, testar de novo e só então fazer merge.
\item Atualizar a cópia local com \codigo{git switch main} e \codigo{git pull}.
\end{enumerate}
\section*{Aplicação no projeto}
O resumo do inventário é uma boa funcionalidade isolada: conta ativos e informa dados úteis sem misturar esse acréscimo com correções de JSON. O histórico do repositório mostra o commit da funcionalidade e sua integração, permitindo explicar quando ela entrou no sistema.
\section*{Checklist demonstrado}
Criar/trocar branch, merge, conflitos, PR e boas práticas colaborativas. A evidência não é só a existência de uma branch: é a PR com alteração compreensível e a \codigo{main} contendo o resultado após a integração.
\cardfonte{https://trello.com/c/kFfXpK39}
\end{document}
